\chapter{Network Level Exploits}
\label{ch:3.1}
\chaptertopic{Topic 3.1}

\chapterquote{The enemy knows the system.}{Claude Shannon (Shannon's Maxim, restating Kerckhoffs's principle)}

\section{Motivation}

Units 1 and 2 covered attacks against people and against physical space. Unit 3 moves into the
network itself -- the plumbing that carries every request between a browser and a server. That
plumbing, TCP/IP, was designed in the 1970s for a small, cooperative research network, not today's
adversarial global internet. It shows: trust is assumed almost everywhere a protocol could
instead demand proof. This chapter works from the ground up -- the protocol stack itself, then
name resolution, then the local-network trust model, then the raw act of listening. Each
later attack in this unit builds on a weakness introduced here.

\begin{apalignment}{Unit 3: Securing Networks}{Topic 3.1 -- Network Vulnerabilities and Attacks}
\begin{apobjectives}
\apobjective{3.1.A}{Analyze Risk}{Name the attacks that recur against networks and say what each
one actually does to the traffic -- capturing frames off a shared medium, rewriting the address
bindings a segment relies on, forging answers that a resolver will cache, forging the source
address in a packet header, or exhausting a switch's finite address table until it fails open.}
\apobjective{3.1.B}{Analyze Risk}{Trace the exploited assumption in each case: which protocol field
is trusted without being verified, which state a receiver keeps on behalf of a stranger, and how an
adversary converts that design decision into stolen data, interrupted service, or destroyed
communication.}
\apobjective{3.1.C}{Analyze Risk}{Assess a network-security decision by what is exposed, what
control actually closes that exposure, and what risk or data loss remains once the decision cannot
be undone.}
\end{apobjectives}
\end{apalignment}

\section{TCP/IP and Topology}

\begin{definition}{Encapsulation}
The process by which each layer of the TCP/IP stack wraps data from the layer above in its own
header, producing a nested structure -- frame, packet, segment, application data -- that is
unwrapped one layer at a time at the destination.
\end{definition}

The stack has four layers, each solving one problem. The \textbf{application layer} is what a user
sees -- HTTP, DNS, SMTP. The \textbf{transport layer} hosts TCP, which breaks data into numbered
segments and delivers them reliably, and UDP, which is faster but guarantees nothing; port numbers
here (443 for HTTPS, 53 for DNS) identify which application should receive the data. The
\textbf{internet layer} is IP, which assigns addresses and routes packets across networks. The
\textbf{network access layer} governs how bits actually move over Ethernet or Wi-Fi, addressing
devices on a local segment by their 48-bit \textbf{MAC address}.

\begin{figure}[h]
\centering
\begin{tikzpicture}[
  layer/.style={draw=kickstartermaroon, thick, align=center, font=\scriptsize},
]
\node[layer, fill=kickstartercream, minimum width=9.6cm, minimum height=0.8cm] (frame) at (0,0) {Frame header (MAC addresses)};
\node[layer, fill=white, minimum width=8.2cm, minimum height=0.8cm] (packet) at (0,-0.9) {Packet header (IP addresses)};
\node[layer, fill=kickstartermaroon!10, minimum width=6.8cm, minimum height=0.8cm] (segment) at (0,-1.8) {Segment header (ports, sequence numbers)};
\node[layer, fill=kickstartermaroon!20, minimum width=5.4cm, minimum height=0.8cm] (data) at (0,-2.7) {Application data (HTTP, DNS, ...)};
\node[font=\scriptsize\sffamily, gray, anchor=west] at (5.1,0) {network access layer};
\node[font=\scriptsize\sffamily, gray, anchor=west] at (4.4,-0.9) {internet layer};
\node[font=\scriptsize\sffamily, gray, anchor=west] at (3.7,-1.8) {transport layer};
\node[font=\scriptsize\sffamily, gray, anchor=west] at (3.0,-2.7) {application layer};
\end{tikzpicture}
\caption{Encapsulation: each layer wraps the layer above it in its own header, producing a nested
structure that is unwrapped one layer at a time at the destination.}
\label{fig:3.1-encapsulation}
\end{figure}

Several structural gaps trace directly back to that 1970s trust assumption. \textbf{IP addresses
are not authenticated} -- a host can claim any source address it likes in a packet header, a
technique called \textbf{IP spoofing}. The receiving host has no built-in way to verify the
claim. Early TCP implementations used \textbf{predictable sequence numbers} in the three-way
handshake (SYN, SYN-ACK, ACK).\srccite{3.1}{1} An attacker who could predict the next number could
inject data into, or hijack, an existing session without ever learning a password. \textbf{ARP operates on pure
trust} on a local segment, a gap explored fully later in this chapter. And \textbf{DNS relies on
caching forgeable replies}, which the next section covers in depth.

\begin{lstlisting}[language=Python, caption={Flagging IP-spoofing indicators in a packet log}]
import ipaddress

def flag_spoofed(packets, expected_subnet):
    """packets: list of dicts like {"src_ip": str, ...}.
       expected_subnet: str CIDR, e.g. "10.0.0.0/24".
       Returns sorted indices of packets whose source IP falls outside
       the expected subnet."""
    network = ipaddress.ip_network(expected_subnet)
    return sorted(i for i, p in enumerate(packets)
                  if ipaddress.ip_address(p["src_ip"]) not in network)
\end{lstlisting}

The network's physical or logical arrangement -- its \textbf{topology} -- shapes the attack surface
just as much as the protocols do. A hub-based topology broadcasts every frame to every port, so any
connected device can passively capture all traffic. Switches fixed this by forwarding frames only
to the correct port. But \textbf{MAC flooding} defeats them by overflowing a switch's address
table until it falls back to broadcasting like a hub. A \textbf{flat network} -- one broadcast
domain, no internal boundaries -- gives every compromised workstation direct visibility to every
server, printer, and camera on it. \textbf{Network segmentation}, dividing the network into
subnets or VLANs with controlled routing between them, limits exactly this blast radius. It is the
subject of Chapter~\ref{ch:3.3}.

\begin{misconception}
\textbf{``IP addresses prove identity.''} An IP address is a routing label, trivially forged; it
is not an authentication token. Real identity verification requires a higher-layer mechanism --
a TLS certificate, an authenticated protocol -- built on top of it.
\end{misconception}

\section{DNS Attacks}

\begin{definition}{DNS cache poisoning}
Tricking a recursive resolver into caching a forged IP-address mapping for a domain, so every
subsequent lookup from any client using that resolver is silently redirected until the cache
entry's TTL expires.
\end{definition}

DNS queries and replies travel over UDP port 53. The protocol is connectionless, with no handshake
and no built-in way for a resolver to verify that an incoming reply actually came from the
authoritative server it queried. A resolver simply accepts the \emph{first} reply that matches three
values from its original query: the destination port it queried from, the exact domain name
requested, and a 16-bit \textbf{Transaction ID}. An attacker who can guess or predict all three,
and deliver a forged reply before the real one arrives, wins a pure race condition. The legitimate
reply arrives late and is discarded.

\begin{example}[The Kaminsky attack]
The classic version of this race forces repeated tries rather than relying on a single lucky
guess. The attacker triggers the resolver to query a non-existent subdomain --
\texttt{1.bank.com}, \texttt{2.bank.com}, and so on -- guaranteeing a fresh query with no cached
entry to short-circuit it. The attacker then floods the resolver with forged replies carrying
guessed Transaction IDs for each one. Because the domain is unique every time, the resolver must
query the authoritative server again and again, giving the attacker many independent attempts to
win the race instead of one. Once a guess lands, the attacker rewrites the authority section of the
reply to redirect the entire \texttt{bank.com} zone, not just the one subdomain
queried.\srccite{3.1}{2}
\end{example}

Modern defenses attack the odds directly. \textbf{Source port randomization} spreads outbound
queries across roughly 64{,}000 ephemeral ports instead of a fixed one. Combined with the
16-bit Transaction ID, this pushes the guessing space from 65{,}536 to over four billion
combinations -- mathematically impractical to brute-force blindly.\srccite{3.1}{2} \textbf{DNSSEC}
adds cryptographic signatures to DNS records, so a resolver discards a spoofed reply for lacking a
valid signature even if the attacker wins the UDP race outright. \textbf{DoH and DoT} encrypt the
link between a client and its resolver, closing local eavesdropping and spoofing on that segment
specifically. They do not protect the resolver-to-authoritative path, though, which still depends
on DNSSEC.

\begin{lstlisting}[language=Python, caption={Detecting a cache-poisoning race in a resolver log}]
def find_poisoning_races(log):
    """log: list of dicts like {"tx_id": int, "domain": str, "answer_ip": str}.
       Returns sorted (tx_id, domain) pairs that received more than one
       distinct answer -- the signature of a poisoning race."""
    answers_seen = {}
    for entry in log:
        key = (entry["tx_id"], entry["domain"])
        answers_seen.setdefault(key, set()).add(entry["answer_ip"])
    return sorted(key for key, answers in answers_seen.items() if len(answers) > 1)
\end{lstlisting}

\begin{misconception}
\textbf{``HTTPS protects against DNS spoofing.''} HTTPS prevents an attacker from reading your
traffic once you are connected, but does nothing to stop DNS from sending you to the wrong server
in the first place. A properly configured browser will reject a forged server's certificate and
warn the user -- but the user still experiences a failed connection or a warning to click through,
not silent protection.
\end{misconception}

\section{Man-in-the-Middle and ARP Spoofing}

\begin{definition}{Man-in-the-middle (MITM)}
An attack in which an adversary secretly positions itself between two communicating parties,
relaying and potentially reading or altering their traffic, while both sides believe they are
talking directly to each other.
\end{definition}

MITM is a category, not one technique. DNS cache poisoning is one way to become a man in the
middle at the naming layer; \textbf{ARP spoofing} is the classic way to become one at the local
network layer. On any Ethernet or Wi-Fi segment, a host must learn a destination's MAC address
before it can send it an IP packet. The \textbf{Address Resolution Protocol (ARP)} solves this
with a broadcast question -- ``who has 192.168.1.5?'' -- and a reply that gets cached. ARP was
designed in 1982 for a small trusted network and has no authentication
whatsoever.\srccite{3.1}{3} Any device can send an ARP reply, even an unsolicited
\textbf{gratuitous ARP} nobody asked for, and every other device updates its cache to match -- no
verification possible.

An attacker already on the local network exploits this directly by sending two forged reply
streams at once. One tells the victim that the gateway's IP now belongs to the attacker's MAC; the
other tells the gateway that the victim's IP now belongs to the attacker's MAC. Both devices update
their tables to point at the attacker instead of each other. If the attacker then simply forwards
traffic on to its real destination after copying it, neither side notices any disruption. The
connection keeps working, just with an invisible extra hop where the attacker can read, log, or
alter everything passing through.

\begin{lstlisting}[language=Python, caption={Flagging ARP-spoofing signatures in an observed-reply stream}]
def flag_arp_spoofing(replies):
    """replies: list of dicts like {"ip": str, "mac": str}, in observed
       order. Returns sorted IPs that were claimed by more than one MAC
       address -- the spoofing signature."""
    macs_by_ip = {}
    for r in replies:
        macs_by_ip.setdefault(r["ip"], set()).add(r["mac"])
    return sorted(ip for ip, macs in macs_by_ip.items() if len(macs) > 1)
\end{lstlisting}

Once positioned as the man in the middle, what an attacker can actually do depends on what
protection the traffic already has. Plain HTTP, unauthenticated DNS, and Telnet are trivially
readable in full. HTTPS-protected traffic cannot be read directly, since TLS encrypts the payload
end to end. But the attacker can still observe which servers are contacted via the TLS SNI field,
attempt an SSL-stripping downgrade to trick the victim into falling back to plain HTTP, or present
a forged certificate that a properly validating browser will reject and warn about.

Several controls answer ARP spoofing specifically. \textbf{Dynamic ARP Inspection (DAI)} is a
switch feature that validates ARP packets against a trusted binding table built from DHCP-snooping
records and drops anything that does not match before it reaches other hosts. \textbf{Static ARP
entries} hard-configure the correct mapping for a high-value host like the default gateway, so its
OS ignores any reply attempting to change it. This is effective, but it does not scale to an
entire large network. \textbf{Network segmentation} limits ARP spoofing structurally, since it only
works within a single broadcast domain. A sensitive host on its own VLAN cannot be spoofed by a
device on a different one at all.

\begin{misconception}
\textbf{``ARP spoofing requires access to the router.''} It requires only that the attacker share
the victim's local broadcast domain -- the same Wi-Fi or switch segment. ARP replies travel
directly between hosts on the LAN, with no router or gateway involvement needed to forge them.
\end{misconception}

\section{Packet Capture and Traffic Sniffing}

\begin{definition}{Packet capture}
Intercepting and logging network traffic as it passes a network interface -- the same technique
tools like Wireshark and tcpdump use for legitimate troubleshooting and that an attacker uses to
eavesdrop on traffic they should not have access to.
\end{definition}

The tool itself is neutral; what makes capture an attack is position (does the attacker have
legitimate access to see this traffic at all?) and content (does the traffic contain something
sensitive in cleartext?). On old hub-based networks, every packet was physically broadcast to
every port, making sniffing trivial for anyone connected. Modern switches send a frame only to the
port holding its destination MAC, which limits a passive listener considerably. Attackers get
around this in one of three ways. The first is ARP spoofing, redirecting traffic through the
attacker's own machine first. The second is abusing \textbf{port mirroring} (a legitimate switch
feature for monitoring tools, copying traffic from one or more ports to a monitoring port) if it is
granted to the wrong party. The third is setting a NIC to \textbf{promiscuous mode}, so it
processes every frame it physically receives rather than only frames addressed to its own MAC --
necessary for a capture tool, but only useful once traffic is actually reaching that NIC by one of
the other two means.

Analysts distinguish two filtering stages. A \textbf{capture filter} is applied before anything is
written to disk -- \texttt{tcp port 80} discards everything else on the wire outright, and packets
never captured cannot later be recovered. A \textbf{display filter} narrows what is shown from a
capture that already contains everything, without discarding any underlying data. Because a capture
filter cannot be undone, analysts commonly capture broadly and filter the display narrowly when
they are unsure exactly what they are looking for.

\begin{lstlisting}[language=Python, caption={Flagging plaintext-credential exposure in a packet log}]
CLEARTEXT_PORTS = {80: "http", 23: "telnet", 21: "ftp"}
CREDENTIAL_MARKERS = ("password=", "authorization: basic", "user ", "pass ")

def flag_credential_leaks(records):
    """records: list of dicts like {"dst_port": int, "payload": str}.
       Returns sorted indices of records that expose a plaintext
       credential based on a cleartext-by-default port and a
       credential-shaped payload marker."""
    flagged = []
    for i, r in enumerate(records):
        if r["dst_port"] in CLEARTEXT_PORTS:
            lowered = r["payload"].lower()
            if any(marker in lowered for marker in CREDENTIAL_MARKERS):
                flagged.append(i)
    return sorted(flagged)
\end{lstlisting}

Any protocol that never encrypts its payload -- plain HTTP, Telnet, FTP, unencrypted POP3/IMAP/SMTP
-- is fully readable by anyone positioned to capture it. That is precisely why the industry has
moved to encrypted equivalents (HTTPS, SSH, FTPS/SFTP). Encryption does not stop an attacker from
capturing the packets; it makes the captured payload useless without the key. That reframing is
also the practical defense here: rather than trying to prevent capture outright, which is often
impossible on a shared network, ensure that anything captured yields nothing of value.

\begin{misconception}
\textbf{``A switched network makes sniffing impossible.''} A switch limits passive sniffing
compared to a hub. But ARP spoofing, port-mirroring misuse, and a compromised network tap all still
let an attacker see traffic they should not, on a fully switched network.
\end{misconception}

\section*{Review Questions}
\begin{reviewquestions}
\item Explain why IP spoofing is possible at all, tracing the explanation back to a specific design
assumption made in TCP/IP's original 1970s context.
\item A resolver log shows two different answer IPs returned for the same transaction ID and
domain. Explain what this pattern indicates and why the resolver accepted the forged one instead
of discarding it.
\item Explain, step by step, how an attacker uses two separate forged ARP replies to position
themselves between a victim and a gateway, and explain why neither party notices a disruption in
service.
\item A network relies entirely on switches (no hubs) and assumes this makes packet sniffing
impossible. Identify two specific techniques that defeat this assumption and explain the
mechanism each one exploits.
\item Explain why DNSSEC stops a cache-poisoning attack even in the scenario where the attacker has
already won the underlying UDP race condition.
\item Using the encapsulation diagram, explain why a device operating only at the network access
layer (a basic Ethernet switch) cannot itself detect IP spoofing, even though the forged field
physically passes through it on every frame.
\item An analyst is unsure exactly what they are looking for and captures traffic broadly with a
minimal capture filter, then narrows what they view with a display filter. Explain why this order
cannot safely be reversed, referencing what a capture filter actually does to unmatched packets.
\end{reviewquestions}

\begin{answerkey}
\item IP spoofing is possible because TCP/IP's original design, standardized in the 1970s for a
small, cooperative research network, never built in a way to verify that a packet's claimed source
address is the address that actually sent it. The source IP field is simply a value the sending
host writes into the header; nothing downstream checks it against the physical origin of the
traffic. Authenticating identity was never part of the protocol's job, so any host can write any
source address it likes with no built-in mechanism to catch the forgery.
\answerlabel{3.1.B}
\item Two different answers for the same transaction ID and domain means at least one of them is
forged -- an attacker raced the legitimate authoritative server and delivered a guessed reply first.
The resolver accepted it because it validates a reply using only three values (destination port,
domain name, and the 16-bit Transaction ID), and caches whichever matching reply arrives first. It
has no way to authenticate the reply's true source. So the forged answer simply won the race, and
the real one, arriving after the cache slot was filled, was discarded on arrival.
\answerlabel{3.1.B}
\item The attacker sends two separate, continuous streams of forged ARP replies. One tells the
victim that the gateway's IP address now maps to the attacker's MAC address; the other tells the
gateway that the victim's IP address now maps to the attacker's MAC address. Because ARP caches any
reply unconditionally, both devices update their tables to point at the attacker instead of each
other. Neither side notices a disruption because the attacker forwards the intercepted traffic on
to its real destination after copying it. The connection keeps working exactly as before, just
with an invisible extra hop where the attacker can read, log, or alter everything passing through.
\answerlabel{3.1.B}
\item MAC flooding overflows a switch's finite address table with forged source addresses until it
can no longer track which MAC sits behind which port. This forces it to fail open and broadcast
every frame to every port like a hub. ARP spoofing achieves the same visibility a different way.
By poisoning the address bindings of a target and a gateway, the attacker's own machine becomes the
path traffic is forwarded through. As a result, the switch is behaving correctly the entire time,
even while it delivers all of the target's traffic to the attacker. Port-mirroring misuse is a
third route, copying legitimate monitoring traffic to a party who should not have received it.
\answerlabel{3.1.A}
\item DNSSEC adds a cryptographic signature to legitimate DNS records that only the holder of the
zone's private key can produce. Winning the UDP race only gets a forged reply accepted \emph{before}
signature verification. A resolver with DNSSEC enabled checks the signature on every answer
regardless of which reply arrived first. A spoofed reply -- lacking a valid signature it has no
way to forge -- is discarded even though it won the underlying timing race outright. The forgery is
defeated at a layer the race condition never touches.
\answerlabel{3.1.B}
\item Encapsulation means each layer only reads and acts on its own header. A switch operating at
the network access layer forwards frames based on MAC addresses, and has no defined role in
inspecting the IP header nested one layer further in. The forged source IP address travels through
the switch completely unexamined, because checking it is not a function the network access layer
performs at all. Detecting IP spoofing requires equipment operating at the internet layer or
above, such as a router or firewall applying ingress filtering.
\answerlabel{3.1.A}
\item A capture filter discards non-matching packets before they are ever written to disk, so any
packet it excludes is gone permanently. There is no later step that can recover data that was
never captured. If the analyst guesses the wrong capture filter while unsure what they are looking
for, the traffic that would have answered their question may already be discarded by the time they
realize the filter was wrong. Capturing broadly first and narrowing only the display afterward
keeps every packet available for re-filtering as many times as needed, at the cost of more storage
and a larger dataset to search through.
\answerlabel{3.1.C}
\end{answerkey}

\section*{Sources}
\begin{enumerate}[label=\lbrack\arabic*\rbrack, leftmargin=2.2em]
\item \srcitem{3.1}{1}Robert T. Morris, ``A Weakness in the 4.2BSD Unix TCP/IP Software,'' Computing
Science Technical Report No.\ 117, AT\&T Bell Laboratories, February 25, 1985 -- the paper that
first formally documented TCP sequence-number prediction as an exploitable weakness.
\item \srcitem{3.1}{2}CERT Coordination Center, Vulnerability Note VU\#800113, ``Multiple DNS
Implementations Vulnerable to Cache Poisoning,'' published July 8, 2008, following Dan Kaminsky's
coordinated disclosure and his August 2008 Black Hat USA presentation of the attack.
\item \srcitem{3.1}{3}D.\ C.\ Plummer, ``An Ethernet Address Resolution Protocol,'' RFC 826,
Internet Engineering Task Force, November 1982.
\end{enumerate}
